% !TEX encoding = UTF-8 Unicode
% !TEX root =  ../Main.tex
% !TEX spellcheck = de_DE

\chapter{Vorgehen und Methodik}\label{cha:Methodik}

Das methodische Vorgehen dieser Arbeit orientiert sich an der gestaltungsorientierten Forschung, auch \glqq{}Design Science\grqq{} genannt. Es beinhaltet die Entwicklung eines Softwareartefakts mithilfe vorhergehender Analyse in Form von Modellen, Methoden oder Systemen. Dieses Softwareartefakt stellt einen Prototyp eines Anwendungsssytems dar und wird nach der Entwicklung evaluiert. \VglZitat{wildeForschungsmethodenWirtschaftsinformatik2007}{281 f.} 
Der Design Science Ansatz erforscht zum einen den Prozess Entwicklung eines Prototypen oder einer Ideen und deren Einfluss auf Prozesse oder Organisationen. Solch ein Prototyp kann beispielsweise eine programmierte Software sein. Der Vorteil darin liegt in der Nähe zur Praxis. \VglZitat{lindnerForschungsdesignsWirtschaftsinformatikEmpfehlungen2020a}{38 ff.} Im vorliegenden Fall handelt es sich dabei um ein digitales Antragssystem zur Erstattung des Semesterticketbeitrags. Der Fokus soll nicht nur auf der technischen Umsetzung liegen, sondern auch eine Analyse des organisatorischen Kontextes mithilfe verschiedener Modelle sowie die Definition von Anforderungen an das System beinhalten. 

Der Design Science Ansatz für den Bereich Informationssysteme wurde besonders bekannt durch Hevners Forschungsarbeit im \glqq{}MIS Quarterly\grqq{}. \NarrativZitat{hevnerDesignScienceInformation2004}{} Darin werden sieben Prinzipien aufgestellt, die als Leitlinie für einequalitative Design-Science-Forschung dienen sollen. Darin wird gefordert, dass die Forschungsergebnisse in Form eines funtkionsfähigen Artefakts vorliegen, ein relevantes Problem adressiert wird, Qualität und Wirksamkeit durch Evaluierung nachgewiesen wird, ein erkennbarer Beitrag zu wissenschaft geleistet wurd, die Forschung auf einem Suchprozess basiert, Iteration ist ein wichtiger Teil davon und dieser Prozess technikorientierten als auch managementorientierten Zielgruppen verständlich erläutert wird. 

Pfeffers entwickelt aus diesen 7 Prinzipien eine Desing-Science-Research-Methologie, die auch in dieser Arbeit Anwendung findet. Diese ist in sechs Prozesstufen unterteilt \VglZitat{peffersDesignScienceResearch2007}{52 ff.}
\begin{enumerate}
    \item \textbf{Problemidentifikation und Motivation:} Das Problem und die Relevanz einer Lösung erläutert. Auf Basis dieser Problemstellung wird in einer späteren Phase das Softwareartefakt als Lösung dieses Problems entwickelt. 
    \item \textbf{Zieldefinition der Lösung:} Die Zielbeschreibung wird rational aus der Problemstellung und dem Wissendarüber abgeleitet, was realisierbar ist. Die Ziele sollten klar definiert sein, quantitativ durch Messgrößen oder qualitativ, beispielsweise anhand einer Beschreibung, wie sich das neue Artefakt verahlten soll.
    \item \textbf{Entwurf und Entwicklung des Artefakts:} In dieser Phase wird das Artefakt erstellt. Diese Erstellung enthält sowohl die Konstruktion, beispielsweise anhand von Modellen als auch die tatsächliche Entwicklung. 
    \item \textbf{Demonstration des Artefakts:} Es folgt die Demonstration der entwickelten Lösung. Das kann anhand eines Experiments, Simulation, Softwaretests oder anderen Methoden erfolgen. Es ist wichtig die Funktionalität des Artefakts darzustellen. 
    \item \textbf{Evaluation:} Es muss kritisch evaluiert werden, wie sehr das entwickelte Artefakt das Problem tatsächlich löst. Das bedeutet, vorher getroffene Annahmen und Zielkriterien aus Phase 2 müssen mit der Realität abgeglichen werden. Nach einem Vergleich kann entschieden werden, ob die Entwurfs- und Konstruktionsphase noch einmal iterativ aufgegriffen werden sollte. 
    \item \textbf{Kommunikation der Ergebnisse:} Dieser Teil hebt hervor, dass die Forschung überzeugend kommuniziert werden muss. Die Erkenntnisse werden dokumentiert und für eine breite Öffentlichkeit (wissenschaftliche als auch praktische Zielgruppen) aufbereitet. 
\end{enumerate}

Phase 1 (Problemidentifikation und Motivation) wird in der Einleitung (Kapitel~\ref{cha:Einleitung}) sowie in der IST-Analyse (Kapitel~\ref{cha:Analyse}) behandelt. Die Zieldefinition (Phase 2) erfolgt in der Anforderungsanalyse ebenfalls in Kapitel~\ref{cha:Analyse}. Der Entwurf (Phase 3) wird anhand des Datenmodells und des Designs verschiedener User-Interface-Komponenten dargestellt, die Entwicklung des Artefakts erfolgt in Kapitel~\ref{cha:Realisierung}. Auf eine Demonstration (Phase 4) wird in dieser Arbeit verzichtet, stattdessen sind die gestalteten User-Interfaces im Anhang dokumentiert, und der Quellcode steht zur freien Erprobung zur Verfügung. Die Evaluation (Phase 5) erfolgt in Kapitel~\ref{cha:Fazit}. Die Kommunikation (Phase 6) der Ergebnisse erfolgt durch die vorliegende Arbeit, die sämtliche Ergebnisse zusammenführt und dokumentiert. 